In this paper, we presented FireRoute, a service-oriented emergency routing and dispatch framework that integrates routing, hydrant availability, environmental sensing, and governance into a unified decision-support workflow. Rather than proposing a new shortest-path algorithm, the contribution lies in demonstrating how heterogeneous operational factors can be systematically combined to produce routing and dispatch decisions that are both actionable and interpretable. The results show that the system responds meaningfully to infrastructure constraints, while also exposing structural fragility that is often overlooked in traditional routing approaches.

From a systems perspective, FireRoute illustrates the value of treating emergency response as a data-driven service pipeline rather than a standalone optimization task. By incorporating explainable risk modeling and evidence-ready governance, the framework supports not only real-time decision-making but also post-incident analysis and accountability. The accompanying reproducibility package further strengthens the contribution by enabling transparent evaluation, scenario replay, and independent verification of results.

Despite these contributions, several limitations remain. The current pilot graph is relatively small and relies on approximate spatial representations. The risk and confidence models are heuristic and not yet calibrated against real-world dispatch data. Environmental factors such as smoke and access constraints are controlled through scenario parameters rather than derived from live sensor fusion. Additionally, blocked edges are modeled using an extreme penalty, which effectively reveals fragility but does not yet incorporate adaptive rerouting strategies or redundancy-aware planning.

Our work will focus on extending the framework toward larger-scale and more realistic deployments. This includes expanding the graph to city-scale networks, integrating real-time sensor streams and historical incident data, and learning parameter weights from operational datasets. Further research will explore adaptive routing strategies that balance feasibility and resilience, as well as multi-objective optimization formulations that jointly consider response time, safety, and resource utilization. In addition, incorporating predictive modeling and uncertainty-aware decision-making could enhance system robustness under incomplete or noisy inputs. Ultimately, these directions aim to transform FireRoute from a controlled pilot system into a scalable and deployable solution for real-world emergency response environments.